\section{Deformations and local geometry}\label{sec:deformation}

We establish the deformation-theoretic tools used in both propagation and
connectedness.  The two points requiring particular care are deformation
theory at a characteristic-two curve and the geometry of several local
conditions imposed simultaneously.  We retain determinant characters
throughout.  In particular, in characteristic two a pseudorepresentation
means a two-dimensional determinant, not merely its trace.

There are three outputs. First, finite trace-and-character spaces
let us compare dimensions before and after forgetting frames and
lines. Second, curve completion gives legitimate continuous local
and global deformation problems with a common coefficient field.
Third, two distinct local geometries serve different later uses:
unrestricted dyadic factors enter propagation, whereas marked scalar
factors enter connectedness.

\subsection{Representation spaces and formal dimensions}

The groups occurring below, namely local Galois groups and the groups
$G_{F,S}$ with $S$ finite and containing the dyadic and infinite places,
satisfy Mazur's finiteness condition $\Phi_2$: for every open
subgroup $H$, the set $\Hom_{\mathrm{cts}}(H,\ZZ/2\ZZ)$ is finite.
We use continuous
deformation functors.  A frame is a specified basis; a local line, when
recorded, is part of the deformation data. This is the framed form
of the deformation framework of \cite{Mazur1989}.

\begin{proposition}[Representation spaces]\label{prop:rep-space}
Let $G$ satisfy $\Phi_2$, let $\overline D$ be a continuous
two-dimensional determinant over the finite field $\kk$, and fix a
continuous lift $\chi$ of its determinant character.  The following
statements hold.
\begin{enumerate}
\item The determinant deformation functor is represented by a complete
Noetherian local $\OO$-algebra $R^{\mathrm{ps}}$.  There is a finite
$R^{\mathrm{ps}}$-algebra $\mathscr E$ carrying its universal Cayley--Hamilton
representation.  Matrix representations of $\mathscr E$ compatible with the
universal determinant form an affine scheme of finite type over
$R^{\mathrm{ps}}$.  Its completion at a closed residual matrix
representation is that representation's framed deformation space.
\item Invariant lines at finitely many closed subgroups can be imposed
by closed incidence conditions in products of projective lines.  On the
absolutely irreducible global locus, existence of such a line with a
specified character is closed in the trace and character parameters.
\item Suppose the subgroups' abelianizations have finitely generated
maximal pro-$2$ quotients.  Adjoining the characters on the specified
lines gives a finite extension of the completed trace-and-determinant
image.
\item At an absolutely irreducible representation over a coefficient
field, determinants classify unframed deformations.  After a splitting
extension, the framed fiber is a conjugacy orbit of dimension three.
\end{enumerate}
These assertions do not require the residual determinant to be
multiplicity-free.
\end{proposition}

\begin{proof}
For representability under $\Phi_2$, we use
\cite[Propositions~3.3 and~3.7, Remark~3.5]{Chenevier2014}, with the
numbering of the cited arXiv version.  Finiteness of the universal
Cayley--Hamilton algebra and the algebraization of its compatible
matrix representations are
\cite[Propositions~3.2 and~3.6, Theorems~3.7 and~3.8]{WangErickson2018}.
Compatibility here includes the characteristic-polynomial equations;
no generalized matrix algebra at a scalar residual point is being
assumed.  Fixing $\chi$ is a base change of these statements.

For clarity, finite type has a direct algebraic explanation.  Choose
finitely many algebra generators for the finite algebra $\mathscr E$, assign a
$2\times2$ matrix to each, and impose the multiplication and
characteristic-polynomial identities.  These are polynomial equations
over the Noetherian base.  Completing at a residual matrix point
recovers continuity and the chosen residual representation, including
its extension class.

Line invariance and the prescribed action on a line are closed
equations in the relative projective line.  Over the absolutely
irreducible locus the Cayley--Hamilton algebra is an Azumaya algebra of
rank four; one may split it faithfully flat locally and then descend
the proper incidence image.  This proves the closedness assertion.

If $g_1,\ldots,g_a$ topologically generate the relevant abelian
character quotient, every recorded character value $\lambda(g_j)$
satisfies
\[
 X^2-T(g_j)X+\chi(g_j)=0.
\]
Modulo the maximal ideal of the trace image, the corresponding
parameters are nilpotent after the chosen residual roots have been
fixed.  The closed fiber of the completed character algebra is
therefore finite.  Complete Nakayama proves finiteness of that algebra
over the trace image.  Finite torsion generators cause no change.

Finally, the absolutely irreducible determinant theorem identifies the
Cayley--Hamilton algebra with a matrix algebra after a splitting
extension \cite[Theorem~2.22 and Corollary~2.23]{Chenevier2014}.  Two compatible representations are conjugate and their
stabilizer consists of scalars.  This gives the final assertion.
\end{proof}

In dimension two the fixed-determinant ring is topologically generated
by traces.  There is no division by two in this statement: for group
elements $g_j$ the determinant of a formal sum is
\[
D\Bigl(\sum_j t_jg_j\Bigr)
=\sum_jt_j^2\chi(g_j)
 +\sum_{j<l}t_jt_l\bigl(T(g_j)T(g_l)-T(g_jg_l)\bigr).
\]
This identity also explains why the determinant character must remain
part of the data in characteristic two.

For a representation chart with recorded local lines, its completed
trace-and-character image is the closed algebra generated by the global
traces and the character values on those lines. Passing to this image
forgets the frame and the positions of the lines, but retains their
characters. Under the abelianization hypothesis of
Proposition~\ref{prop:rep-space}(3), forgetting the characters is then
a finite map and does not change the dimension of the image. The next
lemma bounds the dimension lost in the first passage.

\begin{lemma}[Formal dimension comparison]\label{lem:dimension-comparison}
Let $X$ be one of the finite-type representation spaces of
Proposition~\ref{prop:rep-space}, including any recorded local lines,
over an excellent complete local trace-and-character base $R$.  Let
$x$ be a closed point above the closed point of $R$, and let $P$ be a
prime of $\widehat{\mathcal O}_{X,x}$.  Suppose its generic global
representation is absolutely irreducible.  If $p=P\cap R$, then
\[
\dim\bigl(\widehat{\mathcal O}_{X,x}/P\bigr)
 \leq \dim(R/p)+3+\ell,
\]
where $\ell$ is the number of recorded lines not known to have
distinct characters generically.  A line with distinct characters
contributes zero to this bound.  The analogous statement holds with
any established bound for the relevant generic representation fiber.
\end{lemma}

\begin{proof}
Contract $P$ to a prime $q$ of $\mathcal O_{X,x}$.  There is a
surjection from the completion of $\mathcal O_{X,x}/q$ to
$\widehat{\mathcal O}_{X,x}/P$, and completion preserves dimension
\cite[Tags~02IJ and~07NV]{StacksProject}.
The finite-type dimension formula over the excellent local base
therefore bounds the latter dimension by $\dim(R/p)$ plus the
generic-fiber dimension of the contraction.  Over an algebraic closure
of $\Frac(R/p)$ the representation is unique up to a three-dimensional
conjugacy orbit.  Each unspecified invariant line has dimension at most
one in $\mathbf P^1$; with distinct characters its choices are finite.
This gives the inequality.  Notice that the proof contracts the formal
prime before applying the finite-type formula.  It does not assert
that an arbitrary formal prime itself algebraizes.
\end{proof}

Thus passing from a representation chart to its trace-and-character
image loses at most the three frame dimensions and the dimensions
of the lines whose character has not separated. The next lemma
supplies curves and coefficient points on which these comparisons
can be tested.

\begin{lemma}[Curves and coefficient points]\label{lem:formal-point-density}
Let $A$ be a complete Noetherian local domain with finite residue field.
\begin{enumerate}
\item If $A$ has characteristic two and positive dimension, then its
dimension-one domain quotients have generic points dense in $\Spec A$.
Their normalizations are $\kk'[[t]]$ for finite extensions $\kk'/\kk$.
\item If $A$ is a flat $\OO$-algebra, continuous maps from $A$ to rings
of integers in finite extensions of $\EE$ are dense on $\Spec A$.
\end{enumerate}
In either assertion the points can be chosen outside any prescribed
proper closed subset.
\end{lemma}

\begin{proof}
Use formal Noether normalization, over $\kk$ in the first case and
over $\OO$ in the second; the coefficient and finite normalization
facts used here are \cite[Tags~032A, 032D, 07QW and~07QV]{StacksProject}.  Thus $A$ is finite over a power-series
subring $S$.  Given $0\ne f\in A$, integrality supplies a nonzero
element of $fA\cap S$: use an integral equation for $f$ with nonzero
constant term in the fraction field extension, and clear denominators.

In the first case choose a one-parameter specialization of $S$ on which
that element is nonzero.  Such a specialization exists by choosing
positive weights for the variables and, if necessary, coefficients in
a finite residue extension so that the lowest weighted part does not
vanish.  Its image is finite over $\kk[[t^a]]$ for some $a>0$.
Lying-over gives a one-dimensional prime of $A$ avoiding $f$.
Excellence makes its normalization finite; a complete
equicharacteristic discrete valuation ring with finite residue field
has the asserted form.

In the second case choose sufficiently small integral values for the
variables on which the selected nonzero series does not vanish.
Existence follows successively from the one-variable preparation
theorem.  A prime above this specialization has a finite extension of
$\EE$ as fraction field.  Integrality puts the image of $A$ in its
ring of integers; locality and finiteness make the map continuous.
Again $f$ does not vanish.  Applying this to an element of the ideal of
the excluded closed subset proves density.
\end{proof}

\subsection{Continuous deformations at a curve}

Deformation theory over Laurent-series residue fields, with integral
residual bases and finite coefficient orders, is developed in
\cite[Section~3.5, Proposition~3.34 and Lemmas~3.35--3.37]{BIP}.
We give the completion argument needed here with its global frames
and incidence data. Keeping the Laurent topology makes continuous
cocycles and obstructions the correct ones.

Let $\rho_*:G\to\GL_2(K)$ be a continuous curve representation, where
$K=\kk((t))$.  After a finite extension of constants and of the curve,
choose a stable lattice and write its matrices in $\GL_2(\kk[[t]])$.
Let $\overline\rho$ be their reduction at $t=0$.  A restriction of the
curve to a local group need not generate all of $K$ from its matrix
entries. Introduce a new variable $z$, send it to $t$, and carry this
same variable through every local and global deformation ring. Their
enhanced curve maps then have the common residue field $K$.
The enhanced coefficient ring is
\[
\mathcal L=\widehat{\OO[[z]]_{(\varpi)}}.
\]
It is a complete discrete valuation ring with residue field $K$, and
\[
\mathcal L/\varpi^a=(\OO/\varpi^a)((z)).
\]
Finite-length $\mathcal L$-modules carry their Laurent-series topology.
In particular $K$ is not given the discrete topology.

\begin{lemma}[The curve completion]\label{lem:curve-completion}
Let $R$ be the framed deformation ring of $\overline\rho$, with the
specified determinant and any of the incidence data above.  Let
$\mathfrak p$ be the kernel of the enhanced curve map
$R[[z]]\to\kk[[t]]$.  Then
\[
R_*:=\widehat{(R[[z]])_{\mathfrak p}}
\]
classifies continuous deformations of $\rho_*$ to finite-length local
Artinian $\mathcal L$-algebras with residue $K$.  This identification
commutes with restriction to local groups, with frames and invariant
lines, and with fixing the determinant.  The analogous assertion holds
for determinants; on the absolutely irreducible locus it agrees with
unframed matrix deformation theory.

If a component of the original complete local ring has dimension $e$
and contains the curve, its enhanced completed localization has
dimension $e$.  At a characteristic-zero coefficient point the
corresponding assertion uses the usual completion over its finite
$2$-adic coefficient field, without adjoining $z$; a horizontal
integral component of dimension $e$ then gives a completed local
component of dimension $e-1$.
\end{lemma}

\begin{proof}
We first prove the integral-order assertion underlying the functorial
claim.  Let $C$ be one of the Artinian algebras, with maximal ideal
$\mathfrak n$, $\mathfrak n^N=0$, and $\varpi^a C=0$.  Write
\[
A_0=\im\bigl((\OO/\varpi^a)[[z]]\longrightarrow C\bigr).
\]
The inverse image of $\kk[[t]]$ under $C\to K$ is
$A_0+\mathfrak n$.  A bounded subset of this inverse image lies in a
finite $A_0$-order: subtract coefficientwise lifts of its residual
entries, choose a finite $A_0$-lattice $U\subset\mathfrak n$ containing
the resulting bounded differences, and put
\begin{equation}\label{eq:curve-order}
B=A_0+U+U^2+\cdots+U^{N-1}.
\end{equation}
Powers here mean finite module spans of products.  Nilpotence makes
$B$ an algebra finite over $A_0$.  Its quotient by
$B\cap\mathfrak n$ is $\kk[[t]]$, with nilpotent kernel.  Hence $B$
is local with residue field $\kk$.  The finite local ring $B/zB$ is
Artinian, so a power of $\mathfrak m_B$ is contained in $zB$.
Its maximal-ideal topology is therefore its finite-module topology;
$B$ is complete and compact.

Apply this to the compact set of matrix entries of a continuous lift
of $\rho_*$.  These entries have integral residues because the
residual basis was chosen above.  Their determinants have residual
values in $\kk[[t]]^\times$, and so are units of $B$.  We obtain a
continuous representation in $\GL_2(B)$ reducing to precisely
$\overline\rho$.  Negative powers of $z$ in nilpotent entries belong
to $B\cap\mathfrak n$ and disappear at that closed residual point.
There is no uniform lower valuation bound as $C$ varies, nor is one
needed.  The same construction simultaneously contains any finite
collection of local data, and two such orders admit a common
enlargement.

Universality gives a map $R\to B$.  Adjoin $z$ using its coefficient
value.  Elements outside $\mathfrak p$ have nonzero residue in $K$,
so are units in $C$, while $\mathfrak p^N$ maps to zero after
increasing $N$.  The map factors uniquely through $R_*$.  Conversely,
a local map $R_*\to C$ factors through a finite jet.  Write the
original formal coordinates as integral lifts of their curve values
plus nilpotent increments.  Expansion in the increments has only
finitely many terms at this jet depth; the residual coordinates,
after their constants are removed, lie in $(t)$.  Substitution in
each original formal series consequently converges in the
Laurent-series topology.  It is continuous, and so gives a continuous
lift of the universal representation.  The two constructions are
inverse.  They plainly preserve determinant equations, line
incidences, and restriction maps.  For determinants use their
characteristic coefficients, followed on the absolutely irreducible
locus by the matrix-algebra identification in
Proposition~\ref{prop:rep-space}.

For dimensions, adjoining $z$ increases a component dimension from
$e$ to $e+1$, while the enhanced curve quotient has dimension one:
the map on $\kk[[z]]$ is already onto $\kk[[t]]$.  The dimension
formula for complete equidimensional excellent local rings gives
height $e$ at the enhanced prime.  Localization and completion have
dimension $e$; see also the dimension and catenarity facts in
\cite[Tags~00KW, 032C, 0ECF and~07NV]{StacksProject}.
At a horizontal characteristic-zero coefficient point the integral
quotient again has dimension one, but no new variable has been
adjoined.  Its height and completed dimension are therefore $e-1$.
The deformation functor uses the usual $2$-adic topology.
\end{proof}

The restriction to the inverse image of $\kk[[t]]$ in
\eqref{eq:curve-order} is essential.  For example, the bounded
singleton $\{z^{-1}\}$ cannot lie in a finite order having integral
residual image.

For a representation $V$ over its coefficient field $L$, write
$h^i(G,V)=\dim_L H^i_{\mathrm{cts}}(G,V)$, omitting $G$ when
understood. For global $V$, write
$h^i_v(V)=\dim_L H^i_{\mathrm{cts}}(G_{F_v},V)$.
These are ordinary cohomology dimensions; at real places the
duality pairings use Tate cohomology, as specified below.

\begin{lemma}[Curve-field cohomology]\label{lem:curve-cohomology}
Let $V$ be a finite-dimensional continuous $K$-representation of a
local Galois group or of $G_{F,S}$.  Continuous cohomology is finite
dimensional and satisfies local duality, local and global Euler
characteristics, and Poitou--Tate duality with arbitrary linear local
conditions.  The twist in characteristic two is the residual Tate
twist.  At real places local duality uses Tate cohomology; the
$h^0_v$ appearing in the Euler and Selmer formulas below still
denotes the dimension of ordinary invariants.  In particular, if
$F$ is totally real of degree $d$ and
$\dim_KV=r$, then
\begin{align}
h^0_v(V)-h^1_v(V)+h^2_v(V)&=-r[F_v:\QQ_2]
                       &&(v\mid2),\\
h^0(V)-h^1(V)+h^2(V)&=-rd+\sum_{v\mid\infty}h^0_v(V).
\label{eq:global-euler-curve}
\end{align}
At a real place the local contribution to a Selmer dimension
difference is $\dim L_v-h^0_v(V)$, where
$L_v\subset H^1(F_v,V)$.  The same formulas, with the usual real-place
interpretation, hold for finite $2$-adic coefficient fields.
\end{lemma}

\begin{proof}
Choose a stable $\kk[[t]]$-lattice $T\subset V$.  Continuous cochains
with values in $T$ are the inverse limit of the cochain complexes for
$T/t^nT$.  Finite-level cohomology is finite, so the derived inverse
limits vanish and cohomology commutes with this inverse limit.
The resulting compact cohomology modules are finite over
$\kk[[t]]$: their quotients by $t$ inject into the appropriate
finite-level cohomology group, and compact Nakayama applies.
Every continuous cochain into $V$ has bounded image, since its domain
is compact.  Thus its cochain complex is obtained by inverting $t$,
and the same is true of cohomology.

Apply finite-coefficient local duality and Poitou--Tate to $T/t^nT$
and the dual torsion lattice, then pass to the inverse limit and
invert $t$.  This gives the perfect field-valued pairings, with
Tate cohomology at real places, and the exact global sequence.
Taking lengths in the finite-level Euler
formulas and dividing by $n$ gives the displayed dimensions; bounded
torsion contributes zero in the limit.  The real-place terms pass
through the same calculation and must not be discarded at $2$.
Intersections with lattices and saturation allow any linear local
condition and its annihilator.  The finite-coefficient inputs are
\cite[I, Corollary~2.3, Theorems~2.8, 2.13, 4.10 and~5.1]{Milne2006}.
The identical lattice argument with an $\OO$-lattice gives the
finite-$2$-adic-field version.
\end{proof}

\subsection{Relative presentations}

Curve completion has now identified the deformation functors, and
the preceding cohomology calculation supplies their tangent and
obstruction spaces. We turn these data into a bound on actual
relations cutting out the global ring over its local base. At real
places the displayed invariant dimensions are ordinary \(H^0\),
while the duality argument uses Tate-modified groups.

Put $D=\ad^0(\rho)$ for a fixed-determinant two-dimensional
representation.  Its dimension is three, also in characteristic two.
Its dual there is $\ad(\rho)/\text{scalars}$, not in general $D$
itself.

\begin{proposition}[Relative relation bounds]\label{prop:relative-presentation}
Let $F$ be totally real of degree $d$, let $S$ contain the real and
dyadic places, and let $\rho$ be a representation over a finite
residue field, a curve field as above, or a finite $2$-adic field.
Use one global frame and a fixed determinant.  Let $T\subset S$
consist of all real places and $m$ dyadic places, each isomorphic to
$\QQ_2$.  Over the completed product of the unrestricted local framed
deformation rings at $T$, the global framed ring has a power-series
presentation with $g$ relative variables and $r$ relations satisfying
\begin{equation}\label{eq:relative-one-frame}
g-r\ge 3-6m-h^0\bigl(D^\vee(1)\bigr).
\end{equation}
Local quotient conditions and local line charts may then be imposed
by base change without changing this bound.  This includes replacing
the local base by its reduction.

For the variable-determinant problem with independent frames at the
places of $S$, the relative tangent dimension over the local framed
base is
\begin{equation}\label{eq:relative-independent-frames}
h^1_{\mathrm{Sel}}(\ad\rho)
 +\sum_{v\in S}h^0_v(\ad\rho)-h^0(\ad\rho),
\end{equation}
where the Selmer condition is zero at $S$ and any other prescribed
conditions are included in the Selmer group.  On the absolutely
irreducible locus, forgetting these independent frames has smooth
fiber dimension $4|S|-1$.
\end{proposition}

\begin{proof}
The relative tangent space in the one-frame problem is the kernel of
restriction on $Z^1(D)$.  We spell out the obstruction bound since
it is also used at enhanced curves.  Across a small extension, choose
continuous matrix lifts with the prescribed determinant.
Multiplication errors form a continuous $2$-cocycle.  Given the
specified local lifts, its class belongs to the kernel of restriction
on $H^2(D)$.  If that class vanishes, correct the matrices to a global
representation.  Its differences from the chosen local lifts are
$1$-cocycles, and the second obstruction belongs to the cokernel of
restriction on $Z^1(D)$.  Continuous set sections of finite-length
coefficient modules suffice: coefficientwise digit sections supply
them.  The determinant map on invertible matrices is smooth, so
determinant correction needs no division by two.

Consequently the relation number is at most the sum of these two
obstruction dimensions.  One can verify that this bounds actual
minimal relations as follows.  For a minimal relative presentation
$H\twoheadrightarrow R$ with kernel $I$, test the small extension
\[
H/(\mathfrak m_HI+\mathfrak m_H^j)
 \longrightarrow H/(I+\mathfrak m_H^j)
\]
for $j$ large enough that its kernel is $I/\mathfrak m_HI$.
The two obstruction maps must jointly detect its dual relation
directions. Indeed, quotient the kernel onto an undetected
one-dimensional direction. Vanishing of both obstructions would give
a lift compatible with the local base. Its relative coordinates differ
from the tautological ones by elements of the small kernel, which is
killed by $\mathfrak m_H$. Changing these coordinates therefore alters
a relation only through its residual linear terms in the relative
variables. Those terms vanish by minimality, so the nonzero surviving
relation cannot become zero. This contradicts the existence of the
lift and proves the relation bound for the complete rings.

Poitou--Tate bounds the first obstruction dimension by
\[
h^2(D)-\sum_{v\in T}h^2_v(D)+h^0(D^\vee(1)).
\]
Subtracting it and the $Z^1$ cokernel from the tangent dimension gives
\begin{align*}
g-r
&\ge 3(1-|T|)
 -\bigl(h^0-h^1+h^2\bigr)(D)\\
&\qquad +\sum_{v\in T}
       \bigl(h^0_v-h^1_v+h^2_v\bigr)(D)
 -h^0(D^\vee(1))\\
&=3(1-d-m)+3d-3m-h^0(D^\vee(1)).
\end{align*}
For the last equality, the real terms in
\eqref{eq:global-euler-curve} cancel the real local terms; the dyadic
Euler characteristic is $-3$.  This proves
\eqref{eq:relative-one-frame} in both coefficient characteristics.
Base change carries the same presentation to the extra local
conditions.

For independent frames, first choose a global cohomology class
vanishing at the prescribed local places.  Local choices of
trivialization contribute the local invariant spaces, and a common
global infinitesimal automorphism is quotiented out.  This gives
\eqref{eq:relative-independent-frames}.  If the global representation
is absolutely irreducible, the common automorphisms are exactly the
scalars; the quotient of the $|S|$ frame spaces therefore has
dimension $4|S|-1$.
\end{proof}

\begin{corollary}[A characteristic-zero line chart]
\label{cor:charzero-line-dimension}
For a two-dimensional representation of $G_{\QQ_2}$ over a finite
$2$-adic field with a chosen invariant line and fixed determinant,
every component of its framed line-deformation chart has dimension
at least five.  The odd fixed-determinant real chart has dimension
two.  In particular, over $\QQ$, the relative bound for these two
local charts is $-3$ whenever $h^0(D^\vee(1))=0$.
\end{corollary}

\begin{proof}
Hold the line fixed in an adapted basis.  The coefficient Lie algebra
is the trace-zero upper triangular algebra of dimension two.  The
local Euler formula gives
$\dim Z^1-\dim H^2=4$, so the relation bound gives local dimension
at least four.  Moving the line adds one formal variable.  At a real
place the odd involution has two distinct eigenvalues in
characteristic zero and its conjugacy orbit has dimension two.
The final assertion is \eqref{eq:relative-one-frame} with $d=m=1$.
\end{proof}

\subsection{Small pseudodeformation loci and control of inertia}

\begin{lemma}[Virtually solvable loci]\label{lem:pseudo-bad-loci}
Let $F$ be totally real of degree $d$, let $S$ contain its dyadic and
infinite places, and fix the determinant.  In the characteristic-two
pseudodeformation space, the irreducible dihedral curve points lie in
a closed subset of dimension at most $d$.  Finite-projective-image
curve points give only the residual pseudorepresentation.  If
$F/\QQ$ is abelian, the reducible locus is closed and has dimension
at most one.
\end{lemma}

\begin{proof}
There are only finitely many quadratic extensions of $F$ unramified
outside $S$.  This follows from Hermite--Minkowski, since fixed local
degree bounds the finitely many possible local discriminants.  For
each such extension $F'/F$, use the character deformation space of
$G_{F',S}$ with the product of the character and its conjugate fixed.
Class field theory identifies its free pro-$2$ parameters with a
quotient of the dyadic local units; non-dyadic unit groups and class
groups contribute only finite groups.  The dyadic unit rank over
$F'$ is $2d$, and its minus rank for the quadratic involution is
$d$.  Fixing the product fixes the plus part up to finite groups:
multiplication by two is an isogeny of the relevant free
$\ZZ_2$-modules even though two is not a unit.  Thus the reduced
special character space has dimension at most $d$.

Its map to pseudorepresentations has finite image algebra.  Indeed
the finitely many generating character values satisfy the
characteristic polynomials of the induced traces, and complete
Nakayama applies as in Proposition~\ref{prop:rep-space}.  Its image
is consequently closed and has dimension at most $d$.  Taking the
finite union gives the dihedral bound.  The argument applies to
curve points after finite coefficient and curve extensions, and
Lemma~\ref{lem:formal-point-density} gives the asserted closed-locus
formulation.

If the projective image is finite, each scalar in the representation
has square fixed by its determinant.  The determinant has finite
image in the special fiber, so the scalar image, and hence the whole
image, is finite.  Its traces are algebraic over a finite field.  An
element of $\kk'[[t]]$ algebraic over a finite field is constant;
its constant reduction is already prescribed.  Thus these traces
give exactly the residual pseudorepresentation.

Finally, a reducible determinant is a pair of characters with fixed
product, so one character determines the other.  For an abelian
totally real $F$, Brumer's theorem on Leopoldt's conjecture
\cite{Brumer1967}, in the form proved in
\cite[Theorem~1 and Proposition~8, pp.~9-08--9-09 and 9-14]{Fresnel1969},
gives free pro-$2$ character rank one.
The same finite character-image argument
therefore gives a closed locus of dimension at most one.
\end{proof}

\begin{lemma}[Inertia and solvable base change]
\label{lem:inertia-basechange}
Consider finitely many complete deformation or pseudodeformation
families over finite-residue coefficient rings, and a finite set of
places not above two.  After a fixed finite local extension at each
place, all their field-valued representations have unipotent inertia.
For a characteristic-two curve in such a family whose determinant
has finite image, imposing the residual value of a Frobenius trace
makes the local image finite.

Any further finite local extensions needed to kill these finite
images, to kill the finitely many inertial types, to make the fixed
determinant unramified, or to arrange residue cardinalities
congruent to $1$ modulo $4$, can be realized inside a totally real
solvable base change which splits completely at the dyadic places.
One may take its Galois closure over $\QQ$.  New ramification of the
base field away from the original allowed set creates no new
ramification of the restricted representations.
\end{lemma}

\begin{proof}
Use the finite universal Cayley--Hamilton algebra $\mathscr E$ for each
family.  Its reduction modulo the base maximal ideal is finite.
After killing that finite residual group image locally, the group
acts through the pro-$2$ group $1+\mathfrak m\mathscr E$.  At a place of
residue characteristic $\ell\ne2$, wild inertia and the non-$2$
part of tame inertia therefore act trivially.  Write $\tau$ for the
remaining tame generator and $\phi$ for arithmetic Frobenius, so that
\[
\phi\tau\phi^{-1}=\tau^q.
\]
At any field-valued point the multiset of eigenvalues of $\tau$ is
stable under the $q$th-power map.  Each eigenvalue consequently has
order dividing $q^2-1$.  A further fixed tame extension kills this
semisimple inertia at every point, leaving unipotent inertia.  The
same finite extension works for the finite list of source families.

In characteristic two a unipotent $2\times2$ matrix has square one.
After these extensions the inertia image is therefore finite: it is
generated by the image of the single pro-$2$ tame generator. The
original inertia image is finite as well, since the inertia subgroup
upstairs has finite index in the original one. Now use the Frobenius
chosen in the original local group, whose trace is prescribed by the
curve constraint. Its characteristic polynomial has coefficients in a
finite field, since that trace is the residual constant and the
determinant has finite image. Its eigenvalues have finite order, and
its possible Jordan part has order at most two. This original
Frobenius consequently has finite order. The full local image is an
extension of the finite inertia image by the finite quotient generated
by Frobenius, and is therefore finite.

Every finite Galois group over a nonarchimedean local field is
solvable: wild inertia is an $\ell$-group, and the successive tame
inertia and residue quotients are cyclic.  A composition tower
therefore realizes each required local extension through cyclic
prime-degree steps.  At each step the prime-degree case of the
Grunwald--Wang theorem
\cite[Chapter~VIII, Theorem~2.4]{MilneCFT}
prescribes the finitely many local extensions while imposing
splitting at dyadic and real places.  For degree two this is simply
weak approximation for squareclasses; no higher-power
Grunwald--Wang assertion is needed.  Composita and Galois closure
preserve solvability, total reality, and complete dyadic splitting.
An unramified local extension of even degree makes an odd residue
cardinality $1$ modulo $4$.  The determinant's finite non-dyadic
inertia can be killed at the same time.

Finally, inertia in the enlarged field maps into inertia downstairs.
A representation unramified at a place outside the original allowed
set remains unramified after restriction, even if the field extension
itself ramifies there.  This proves the last assertion.
\end{proof}

The unipotent assertion may be carried to closures of restricted
pseudorepresentation images.  On the globally absolutely irreducible
locus this can be checked by characteristic polynomials and by
traces against a matrix-algebra basis of group elements.  For
example, the equations
\[
\tr\bigl((\rho(\tau)-1)^2\rho(g)\bigr)=0
\]
for such a basis force $(\rho(\tau)-1)^2=0$ by the perfect trace
pairing on $M_2$.  Thus these are closed identities on the locus
where the corresponding trace Gram matrix is invertible.  This
justifies the set-theoretic use of the uniform inertia assertion on
the source-image closures below.

\subsection{The local product used in propagation}

We now define the local model independently of any patched global
ring.  Fix $d$ real factors and $d$ dyadic factors, the latter all
over $\QQ_2$.  At each use the prescribed fixed determinant, odd at
the real places.  Start with integral residual bases over a common
finite field, take the completed product of the framed local rings,
and form its enhanced completion at any chosen curve as in
Lemma~\ref{lem:curve-completion}.  Denote the resulting
$\mathcal L$-algebra by $A^{\mathrm{rd}}$.

We enlarge $\OO$ once so that it contains $i$ with $i^2=-1$.
Choose $u$ odd residue cardinalities $q_j\equiv1\pmod4$, and $g\ge0$.
Put
\[
A=A^{\mathrm{rd}}[[X_1,\ldots,X_g]].
\]
For $\eta\in\{1,i\}$ define $B_\eta$ by adjoining at each of the $u$
places two formal invertible matrices $\tau_j,\phi_j$, both reducing
to the identity at the curve, and imposing
\begin{equation}\label{eq:tame-model}
\tr\tau_j=\eta+\eta^{-1},\qquad
\det\tau_j=1,\qquad
\phi_j\tau_j\phi_j^{-1}=\tau_j^{q_j}.
\end{equation}
There is no condition on $\det\phi_j$.  This is the raw completed
matrix quotient: we do not silently replace it by its reduction.
At characteristic-zero field points the $i$-condition gives the two
distinct inertia eigenvalues $i,-i$, whereas the $1$-condition makes
inertia unipotent. Both $i$ and $-i$ reduce to one in characteristic
two. The two conditions thus have identical special-fiber equations
but different characteristic-zero component geometry, described below.

\begin{proposition}[Local component geometry]\label{prop:local-geometry}
The ring $A^{\mathrm{rd}}$ is a flat complete-intersection normal
domain over $\mathcal L$, of dimension $1+8d$, and its special fiber
is a normal domain of dimension $8d$.  Set
\[
e=1+8d+4u+g.
\]
Then the following statements hold for the model
\eqref{eq:tame-model}.
\begin{enumerate}
\item $\Spec B_i$ is irreducible of dimension $e$, and its unique
irreducible component is horizontal.
\item The two rings have a canonical common special fiber, compatible
with their universal matrices and the variables $X_1,\ldots,X_g$.
Its irreducible components have dimension $e-1$ and are indexed by a
choice, at each tame place, of the unramified component or the closure
of the nontrivial-unipotent component.
\item The irreducible components of $\Spec B_1$ are horizontal and
have dimension $e$.  They are indexed by the same choices.  The
specialization of each has precisely its corresponding special-fiber
component as underlying support.  In particular the generic point
of a special-fiber component lies on exactly one component of
$\Spec B_1$.
\end{enumerate}
The assertions remain valid after finite extension of the constants
and curve field and the resulting enhanced completion.  No condition
on the dyadic curve representations is required.
\end{proposition}

\begin{proof}
\emph{The unrestricted real and dyadic base.}
We first establish the assertion about $A^{\mathrm{rd}}$.  An odd
real framed ring is given by trace zero and determinant $-1$.
These equations are equivalent to the involution relation with that
determinant, by Cayley--Hamilton and invertibility.  At the scalar
residual point it has the form
\[
\OO[[a,b,c]]/(a^2+2a+bc).
\]
Its special fiber $a^2+bc=0$ is a geometrically normal surface: it is
an integral hypersurface whose singular locus has codimension two
\cite[Tag~031S]{StacksProject}.
At a non-scalar residual involution one obtains a smooth chart of
the same relative dimension two.

For a dyadic factor, the fixed-determinant framed deformation theorem
\cite[Theorem~1.5]{BIP} gives a flat complete-intersection normal
domain of dimension seven and a normal special-fiber domain of
dimension six.  The theorem applies at $2$, including scalar and
reducible residual representations; the prime restriction on a
different density theorem is not being used.  The same statement
holds after each finite coefficient extension.

The completed tensor product of these special-fiber rings over the
finite, hence perfect, field is a normal domain.  Here we use the
completed tensor-product normality theorem
\cite[Theorem~3.3 and Proposition~3.1]{Tabaa2020}; the residue tensor
product is the field itself, so the product is local and connected.
Its dimension is $8d$.  The mixed-characteristic completed product
is a flat complete intersection: concatenate presentations, first
checking the regular sequences modulo $\varpi$, and then lift them
together with the nonzerodivisor $\varpi$.  Its special fiber is the
normal domain just obtained, and the regular-element normality
criterion makes the product normal.  It is a domain as well: the
lowest nonzero $\varpi$-adic terms of a putative product of two
nonzero elements could not multiply to zero in the special fiber.

Adjoining $z$, localizing at the enhanced curve, and completing
preserve normality by excellence
\cite[Tags~07QW and~0C23]{StacksProject}.  They preserve the asserted
dimensions by Lemma~\ref{lem:curve-completion}, and the same applies
to the special fiber.  The element $\varpi$ remains a
nonzerodivisor, so the enhanced ring is flat over the discrete
valuation ring $\mathcal L$.  This proves the assertion about
$A^{\mathrm{rd}}$ and also shows why no special dyadic hypothesis is
needed.  The construction can be repeated after any of the finite
extensions in the statement.

\emph{The integral incidence for the \(i\)-type.}
We prove the tame assertions using proper incidences, keeping their
images distinct from the incidences themselves.  In the $i$-case,
Cayley--Hamilton gives $\tau_j^2=-1$, hence
$\tau_j^{q_j}=\tau_j$.  Adjoin a $\phi_j$-stable $i$-eigenline for
each $\tau_j$.  This gives a closed proper incidence $Y_i$ in a
product of projective lines over $\Spec B_i$.  In an adapted basis
on a slope chart, its matrices are
\[
\tau_j=\begin{pmatrix}i&n_j\\0&-i\end{pmatrix},\qquad
\phi_j=\begin{pmatrix}1+d_j+x_j&h_j\\0&1+d_j\end{pmatrix},
\]
and their sole remaining equation is
\begin{equation}\label{eq:i-incidence}
x_jn_j=2i h_j.
\end{equation}
The line slopes are free chart coordinates.

Here is a completion argument which also works for all places at
once.  On a product of affine slope charts put
\[
C=\widehat{A[s_1,\ldots,s_u]}^{\,\mathfrak m_A}.
\]
This is completion along the entire closed fiber, not at a selected
slope origin.  Restricted coefficient series embed it into
$A[[s_1,\ldots,s_u]]$, so $C$ is a domain.  Exactness of completion
identifies $C/\varpi C$ with the corresponding restricted-series
ring over $A/\varpi A$, also a domain.  Completion of the incidence
chart along that fiber is
\begin{equation}\label{eq:i-incidence-completion}
C[[x_j,n_j,h_j,d_j:1\le j\le u]]/
                  (x_jn_j-2ih_j)_{j=1}^u.
\end{equation}
Reduce each monomial by replacing paired factors $x_jn_j$ by
$2ih_j$ until they never occur together.  Assign weights $1,1,2$
to $x_j,n_j,h_j$.  Reduction preserves weight, so it converges in
the power-series topology and yields a normal form.  The difference
between a monomial and its reduction is a sum of the relations
\eqref{eq:i-incidence} times series whose weights are lower by two;
these series converge as well.  The
substitution
\[
h_j\longmapsto x_jn_j/(2i)
\]
embeds these normal forms into
$\Frac(C)[[x_j,n_j,d_j:1\le j\le u]]$.  Indeed the two output exponents
recover the exponent of $h_j$ as their minimum, and each output
coefficient receives only finitely many input terms.  This proves
uniqueness of the normal form and shows that
\eqref{eq:i-incidence-completion} is a domain.  This argument needs
no uniform bound on denominators.  It also shows that the successive
relations are nonzerodivisors and gives dimension
$\dim A+4u=e$ on the incidence.

We explain why chart completions suffice.  At every point of the
closed fiber, localization of the chart maps faithfully flatly to a
localization of its completed chart.  The latter is a domain, so the
incidence is locally integral there.  Every irreducible component of
a proper scheme over a local base meets the closed fiber; so does
any nonempty closed intersection of components.  Local integrality
therefore makes the components disjoint.  The closed fiber is
$(\mathbf P^1_K)^u$, which is connected, hence there is only one
component.  A nonempty nilradical support would also meet the closed
fiber, so the incidence is reduced.  Thus $Y_i$ is integral.

After inverting two the eigenline is uniquely given by the idempotent
$(\tau_j+i)/(2i)$, and the incidence projection is an isomorphism.
In characteristic two, commuting matrices have a common eigenline
after algebraic extension, so the proper projection covers every
geometric point as well.  Its image is therefore all of
$\Spec B_i$ set-theoretically.  The integral incidence is horizontal
and generically finite over that image.  This proves assertion (1).
It does not assert that the raw ring $B_i$ is reduced.

\emph{The horizontal unipotent patterns.}
For $\eta=1$, write $N_j=\tau_j-1$.  The equations give
$N_j^2=0$ and $\phi_jN_j=q_jN_j\phi_j$.  At a field point either
$N_j=0$, or its kernel and image are the same line.  The first
alternative is the unramified factor, with $\phi_j$ unrestricted.
For the closure of the second alternative use the proper line
incidence with adapted matrices
\begin{equation}\label{eq:unipotent-incidence}
\tau_j=\begin{pmatrix}1&n_j\\0&1\end{pmatrix},\qquad
\phi_j=\begin{pmatrix}q_j(1+d_j)&h_j\\0&1+d_j\end{pmatrix}.
\end{equation}
At a nonzero $N_j$ the relation forces exactly this diagonal ratio.
Its chart has three formal matrix coordinates and one slope.

Choose either alternative at each place.  The resulting proper
incidence has free completed slope charts over $A$, so the same
closed-fiber argument proves that it is integral, horizontal, and
of dimension $e$.  Its projection is an isomorphism where every
chosen ramified $N_j$ is nonzero, because the line is then
$\im N_j=\ker N_j$.  Each image is consequently an integral
horizontal closed subset of dimension $e$.  These finitely many
images cover all geometric points of $\Spec B_1$.

\emph{Specialization and uniqueness of the containing component.}
The dimension calculation alone is not enough for support transfer:
we must know exactly which horizontal component contains each generic
special-fiber point. It remains to identify their specializations and
distinguish them.
The two raw presentations reduce to the same one modulo $\varpi$,
since $i\equiv1$ and $q_j$ is odd.  At one place write
\[
\tau=\begin{pmatrix}1+a&b\\c&1+a\end{pmatrix},\qquad
\phi=\begin{pmatrix}1+d&e'\\f&1+g'\end{pmatrix}.
\]
The common special matrix factor is
\begin{equation}\label{eq:tame-special-ring}
K[[a,b,c,d,e',f,g']]/
 (a^2+bc,\ b(d+g'),\ c(d+g'),\ bf+ce').
\end{equation}
Its two reduced components are
\[
U:\ a=b=c=0,
\qquad
R:\ g'=d,\quad a^2+bc=bf+ce'=0.
\]
Both have dimension four and are geometrically integral.  For $R$,
prove this first in the homogeneous polynomial quotient.  On $b\ne0$
one eliminates $c=a^2/b$ and $f=a^2e'/b^2$, obtaining a domain.
The locus $b=0$ has dimension three, whereas every component of the
two-equation quotient has dimension at least four.  There is thus
only one component, generically reduced.  The quotient is a
complete intersection and has no embedded primes, so it is a
domain.  Completion is the product of its finite homogeneous pieces;
lowest-degree nonzero terms show that it remains a domain.  The
same proof works over every field extension.  This argument is
made before completion; it does not substitute into arbitrary
formal series after inverting $b$.

For several places over $A/\varpi A$, the polynomial products of
the chosen $U$ and $R$ factors are flat over that base and have
geometrically integral generic fiber.  They are domains, and the
same homogeneous completion argument applies, completing the base
coefficients and then the positive-degree variables.  Each pattern
therefore gives an integral special-fiber component of dimension
$e-1$.  Reduction of its proper incidence is integral and maps
onto precisely this component.  Proper surjectivity onto the image
is preserved by base change, so this is the entire underlying
specialization of the corresponding horizontal image.

The patterns are distinct.  On the generic ramified special piece
$\tau_j\ne1$, whereas every unramified piece has $\tau_j=1$.
On the generic unramified piece $\tr\phi_j\ne0$, whereas every
ramified special piece has $\tr\phi_j=0$.  Thus no horizontal image
contains the generic special point of a different pattern.  The
images already cover the raw spectrum, proving assertions (2) and
(3), including uniqueness of the horizontal component through
each special-fiber generic point.
\end{proof}

\begin{remark}
The common raw special fiber need not be reduced.  In
\eqref{eq:tame-special-ring}, the element $a(d+g')$ can be nonzero
with square zero.  The proper incidence is also not a normalization:
over $\tau=\phi=1$ its fiber contains $\mathbf P^1$.  Our conclusions
concern precisely the spectra, components, and universal matrix
actions required for support propagation, and do not make either
of these stronger assertions.
\end{remark}

\subsection{Several marked scalar places}

The marked local condition used in connectedness is different from
the unrestricted dyadic condition.  At a scalar residual dyadic
representation, record an invariant line and retain no other local
restriction.  We need the reduced characteristic-two local base,
jointly with the real factors, to be integral. Connectedness will use
this entire completed product as a domain, so integrality of each
individual dyadic factor does not suffice.

\begin{lemma}[Simultaneous marked factors]\label{lem:marked-local-base}
Over a finite characteristic-two field, take $d$ fixed-determinant
real framed factors and $m$ dyadic framed factors whose residual
matrices are scalar, with one invariant line recorded at each
dyadic factor.  The determinant is fixed, and after the residual
scalar twist its dyadic reduction is one.  The reduction of their
completed local product is an integral complete-intersection base
of dimension
\[
2d+5m.
\]
It remains integral after a finite extension of the residue field.
Equivalently, one may first kill the nilpotent torsion parameters in
the diagonal character rings; the resulting completed product is
already reduced and integral and has the same spectrum as the
original product.
\end{lemma}

\begin{proof}
Move a marked line in its one-dimensional formal chart and choose
the corresponding adapted basis.  After the scalar twist, the
diagonal entries are characters $\lambda,\lambda^{-1}$ congruent
to one.  Their reduced character ring is
\[
C_0=\kk[[x,y]].
\]
Indeed $\QQ_2^\times$ has maximal pro-$2$ quotient
$\ZZ/2\ZZ\times\ZZ_2^2$, and the first factor contributes only a
nilpotent character parameter in characteristic two.

The maximal pro-$2$ Galois group of $\QQ_2$ has three generators
and one relation. The odd-degree dyadic case was established by
Serre \cite[Sections~4.4 and~8]{Serre1964Demushkin}; in Labute's
classification, the $q=2$, odd-local-degree case gives generators
with relation
$g_1^2g_2^4[g_2,g_3]=1$
\cite[\S5, Theorem~8]{Demushkin}. If $b(g)$ is the upper-right entry
of the triangular matrix with diagonal entries
$\lambda(g),\lambda(g)^{-1}$, put $c(g)=\lambda(g)b(g)$. Matrix
multiplication gives the left-cocycle identity
\[
c(gh)=c(g)+\lambda(g)^2c(h).
\]
Its values $u_j=c(g_j)$ satisfy one homogeneous linear equation
\begin{equation}\label{eq:marked-cocycle}
a_1(x,y)u_1+a_2(x,y)u_2+a_3(x,y)u_3=0.
\end{equation}
All these deformations factor through the maximal pro-$2$ quotient:
their reduction is the identity and their finite coefficient
quotients have pro-$2$ congruence image.  Thus this presentation
does describe the complete triangular deformation ring after the
diagonal reduction, rather than just a subfunctor of it.

Write $B=1+x=\lambda(g_2)$ and $C=1+y=\lambda(g_3)$.
The reduced diagonal relation gives $\lambda(g_1)=B^{-2}$.
Using $[g_2,g_3]=g_2^{-1}g_3^{-1}g_2g_3$ and this left-cocycle
identity, direct substitution gives
\[
 (a_1,a_2,a_3)=
 \left(\frac{x^4}{B^4},\;
       \frac{x^6}{B^8}+\frac{y^2}{B^2C^2},\;
       \frac{x^2}{B^2C^2}\right).
\]
Multiply \eqref{eq:marked-cocycle} by the unit $B^2C^2$ and set
\[
 v=u_3+\frac{x^2C^2}{B^2}u_1+\frac{x^4C^2}{B^6}u_2.
\]
This triangular formal change of variables transforms the relation
into $x^2v+y^2u_2=0$.  Thus the coefficient ideal is exactly
$(x^2,y^2)$ and its radical is $(x,y)$; the ideal itself is not
radical.

Let $A_\infty$ be the completed real special-fiber product.  The
real calculation in Proposition~\ref{prop:local-geometry} makes it
an excellent Cohen--Macaulay domain of dimension $2d$.  Adjoin the
independent pairs $x_j,y_j$ and the line variables $s_j$, obtaining
\[
A_0=A_\infty[[x_j,y_j,s_j:1\le j\le m]].
\]
Before completing the upper entries, form the graded algebra
\begin{equation}\label{eq:marked-polynomial-product}
P=A_0[u_{jr}:1\le j\le m,\ 1\le r\le3]/(f_1,\ldots,f_m),
\end{equation}
where each $f_j$ is the corresponding relation
\eqref{eq:marked-cocycle}.  The coefficients of $f_j$ involve only
$x_j,y_j$ and have common-zero locus $(x_j,y_j)$.

On a base stratum where exactly $r$ coefficient rows vanish, the
base loses $2r$ dimensions and the upper-entry fiber gains only
$r$ dimensions.  Thus the corresponding locus in $\Spec P$ has
dimension at most
\[
\dim A_0+2m-r.
\]
Every minimal prime of the homogeneous relation ideal is homogeneous
and is contained in $(\mathfrak m_{A_0},u_{jr})$.  In that excellent
equidimensional local ambient ring, the height theorem for the
$m$ equations makes every component have dimension at least
$\dim A_0+2m$.  No component can therefore lie over a nonempty
rank-drop stratum.  Where all coefficient rows are nonzero, the
quotient is a vector bundle of rank $2m$ over an integral base.
There is consequently only one component, generically reduced.

The same argument for every initial subset of the equations shows
that their heights are the expected ones.  They form a regular
sequence in the Cohen--Macaulay ambient ring.  The quotient has no
embedded associated primes, and generic reducedness now proves
that $P$ is a domain.

Finally, each homogeneous piece $P_n$ is a finite, torsion-free
$A_0$-module.  Over $F_0=\Frac(A_0)$ the $m$ independent linear
relations eliminate one upper-entry variable in each triple, so
\[
\prod_{n\ge0}P_n
 \lhook\joinrel\longrightarrow
\prod_{n\ge0}(P_n\otimes_{A_0}F_0)
 \simeq F_0[[v_1,\ldots,v_{2m}]]
\]
is an embedding of rings.  Since $A_0$ and its finite modules are
complete, the left side is precisely the completed local product
in question.  It is a domain, of dimension
$\dim A_0+2m=2d+5m$.  It is also a complete intersection by the
regular sequence already exhibited.

Killing the original diagonal torsion parameters changes no prime
ideal: each is nilpotent, and the ideal they generate is
nilpotent.  The obtained quotient is already reduced, proving the
assertion about the original reduction.  The real normality and
the cocycle and dimension arguments hold after every finite
constant extension, proving the final assertion.
\end{proof}

The two local outputs will be used separately. Propagation uses the
unrestricted real--dyadic base and the correspondence between the
horizontal tame components and their common special fiber.
Connectedness uses the integral marked base of dimension \(2d+5m\).
The latter statement includes no normality assertion, and neither
construction replaces the raw tame rings by their incidences.
